\subsection{Control, Feedback, and Operator Viewpoints}

The Koopman-style material belongs to the most advanced end of this chapter.
Its business value is that it keeps feedback in view. Behavioural systems are
not open-loop. They observe, act, receive error, and act again.

\begin{discussion}
Operator-based viewpoints offer a possible formal upgrade path because they aim
to represent nonlinear controlled dynamics through transformed linear objects.
That does not simplify the human problem into something solved. What it does is
provide a rigorous long-run direction for later appendices: how might one study
controlled nonlinear transition systems without pretending the control law is
absent?
\end{discussion}
